% !TEX root = ../main.tex
\section{Discussion and Improvements}

\textbf{Correctness, completeness and optimality.} The grid is finite and every graph search keeps a
closed set, so UCS and \astar{} are complete: each returns a path if one exists and reports failure
otherwise (test run 3). With unit step costs, UCS is optimal. \astar{} with consistent Manhattan
distance is optimal and never needs to reopen a cell. Every run in Section~5 confirmed equal costs.
The backtracking CSP solver is complete: forward checking only removes values that cannot appear in
any solution extending the current assignment. So if a schedule exists it is found, and otherwise
\code{None} is returned. The brute-force comparison in the test suite confirms this on small
instances.

\textbf{Performance.} \astar{} never did worse than UCS in nodes expanded. Its saving ranges from
none on open, symmetric layouts to linear-versus-quadratic on corridor-like routes. The CSP
instances are tiny; with MRV $+$ degree, LCV and forward checking, the sample and our schedule are
solved without backtracking.

\textbf{Limitations.}
\begin{itemize}
  \item The visiting order is fixed by the slot order. The robot's total route is therefore optimal
    \emph{for that schedule}, not the shortest tour over all feasible schedules.
  \item One unreachable task stops the whole run (\dec{15}). Later reachable tasks are not
    attempted, and the robot is not returned to the charger.
  \item Movement is 4-directional with unit cost. There is no battery, time-per-move, or
    dynamic-obstacle model.
  \item Paths are recomputed for every leg and every algorithm. There is no caching of distances.
\end{itemize}

\textbf{Improvements.}
\begin{itemize}
  \item Choose among feasible schedules by route length: compute pairwise distances between all
    stops with search, then minimise the tour cost subject to the CSP constraints (for example,
    branch and bound over the CSP solutions).
  \item On an unreachable task, skip it, continue with the remaining reachable tasks, and return to
    the charger, then report partial success.
  \item Use stronger propagation (AC-3) or detect precedence cycles before search, for larger
    schedules.
  \item Tie-break \astar{} on larger $g$ among equal $f$, which reduces expansions on open grids
    where many cells tie (the corner case in Table~2).
\end{itemize}
